\section{The M-CTX Framework}\label{sec:overview}

\begin{figure*}[t]
\centering
\includegraphics[width=0.95\textwidth]{fig_pipeline.pdf}
\caption{The M-CTX architecture. An anchor batch is dispatched in
parallel to three composable index families---BR-LZ for OSM range queries
(Section~\ref{sec:brlz}), a linear-time SDF engine (Section~\ref{sec:sdf}),
and a streaming B$^{x}$-tree for AIS $k$NN (Section~\ref{sec:neighbor}).
An optional Joint Context Index (Appendix~\ref{app:jcx}) shares the
per-anchor spatial localisation across all three. The emitted context
tensors are numerically identical to the reference pipeline's, so any
downstream model runs unchanged.}\label{fig:pipeline}
\end{figure*}

M-CTX realises the workload $\mathcal{C}$ of Section~\ref{sec:problem} as a
composition of three independent, individually optimal indices behind a
single API. The \textbf{OSM-Index} answers $R_F$ via three interchangeable
back-ends---a pure-NumPy STR-tree, the BR-LZ learned index, and a
production R$^{*}$-tree---all exact. The \textbf{SDF engine} computes
$D$ in linear time. The \textbf{Neighbor-Index} answers $K_S$ with an
incremental B$^{x}$-tree. Because $R_F$ depends only on location, $D$ on
the clipped patch, and $K_S$ on spatiotemporal location, the three are
independent across anchors and the framework vectorises over the batch
dimension and partitions across workers (Section~\ref{ssec:concurrent}).

The framework is invoked through one entry point,
\begin{verbatim}
mctx.build_context(
  anchors, osm_index,
  sdf_engine, neighbor_index)
\end{verbatim}
returning a dictionary whose schema matches the reference context output
\emph{exactly}. This drop-in property is a design invariant, not an
afterthought: it is what lets us substitute every stage with an indexed
implementation and still guarantee, in Section~\ref{sec:downstream}, that
the downstream model's predictions are unchanged to the bit. The
acceleration is therefore free of any accuracy trade-off, which is the
property a production training pipeline requires before it will adopt a
new context backend.

Because the three sub-queries are independent across anchors, M-CTX is
embarrassingly parallel: the reference implementation vectorises each
stage across the anchor-batch dimension with NumPy/PyTorch and partitions
anchors across worker processes with a \texttt{mapPartitions}-style
driver. A single broadcast-loaded copy of each index serves all workers,
so the framework scales with deployment width without re-building state
(Section~\ref{ssec:concurrent}). The same composition runs unchanged on a
single laptop core and on a multi-socket server, which is the portability a
training pipeline embedded in heterogeneous infrastructure requires.
